\scanpage{82}
\mspage{81}
the $E_i$ which have \unsure{isopyptic|isoclinic|isopyntic|isopypic|isotypic} components $\handwritten{\lambda}$ suc\rlap{l}h that $0\handwritten{<\lambda\leq\leq}1$ are ex\rlap{q}actly the Dieudonné modules of $p$-divisible formal groups over $k$ (the category of these Dieudonné modules \deleted{corresponds to}\inserted{is opposite to}\handwritten{\begin{tikzpicture}[baseline=0cm,x=1cm,y=1cm]\draw[thin] (0,-.18)..controls (-.08,.1) and (.05,.18)..(.22,.18)..controls (.8,.26) and (1.3,.2)..(2.3,.2)..controls (2.4499999999999997,.2) and (2.48,.3)..(2.42,.36);\end{tikzpicture}} the category of these groups, modulo isogeny). On the other hand, an F-isocrystal has slope amplitude contained in $\handwritten{\rlap{\begin{tikzpicture}[baseline=0cm,x=1cm,y=1cm]\draw[gray,thin] (.13,.48)..controls (-.1,.4) and (-.08,-.13)..(.14,-.14);\end{tikzpicture}}[}0,1\handwritten{\rlap{\begin{tikzpicture}[baseline=0cm,x=1cm,y=1cm]\draw[gray,thin] (-.05,.47)--(.1,.49);\draw[gray,thin] (-.02,-.05)..controls (.02,-.17) and (.1,-.17)..(.13,-.13);\end{tikzpicture}}]}$ (closed interval) if and only if it i\rlap{d}s the Dieudonné module of a Barsotti-Tate group over $k$ (the category of F-isocrystals e\unsure{\rlap{b}n|\rlap{h}n|n}\rlap{i}v\rlap{v}isaged being \inserted{anti}\handwritten{\begin{tikzpicture}[baseline=0cm,x=1cm,y=1cm]\draw[thin] (0,-.18)..controls (-.08,.1) and (.05,.18)..(.22,.18)..controls (.8,.26) and (1.3,.2)..(0.75,.2)..controls (0.9,.2) and (0.9299999999999999,.3)..(0.87,.36);\end{tikzpicture}}equivalent to the category of Barsotti-Tate groups over $k$ modulo isogeny ..): the unit F-isocrystal co\handwritten{rr}esponds to the ind-étale Barsotti-Tate group $\handwritten{\underline}{Q}_p/\handwritten{\underline}{Z}_p$ over $k$.

Still assuming $k$ perfect, th\rlap{q}e condition \deleted{of the exists} that $(E,F_E)$ be effective \rlap{a}and that there exists an \inserted{effective}\handwritten{\begin{tikzpicture}[baseline=0cm,x=1cm,y=1cm]\draw[thin] (0,-.18)..controls (-.08,.1) and (.05,.18)..(.22,.18)..controls (.8,.26) and (1.3,.2)..(1.6,.2)..controls (1.75,.2) and (1.78,.3)..(1.7200000000000002,.36);\end{tikzpicture}} $V_i$ such that $FV_i=V_iF=p^i\operatorname{id}$ (cf number\textsuperscript{o}2) means that $E$ has slope amplitude contained in $\handwritten{\rlap{\begin{tikzpicture}[baseline=0cm,x=1cm,y=1cm]\draw[gray,thin] (.13,.48)..controls (-.1,.4) and (-.08,-.13)..(.14,-.14);\end{tikzpicture}}[}0,i\handwritten{\rlap{\begin{tikzpicture}[baseline=0cm,x=1cm,y=1cm]\draw[gray,thin] (-.05,.47)--(.1,.49);\draw[gray,thin] (-.02,-.05)..controls (.02,-.17) and (.1,-.17)..(.13,-.13);\end{tikzpicture}}]}$.

10. \handwritten{\uline}{Case} $k$ \handwritten{\uline}{finite}. Suppose $k=\handwritten{\underline}{F}_q^{\deleted{\handwritten{\unsure{*|p|\overline{\vphantom{F}}}}}}$, where $q=p^{\deleted{a}\handwritten{\nu}}$. By virtue of 6.1, \deleted{its} the \unsure{\rlap{e}g|g}erbe $\handwritten{\underline}{G}(\handwritten{\underline}{F}_q)$ is known, it is the obvious subgerbe of $\handwritten{\underline}{G}(\handwritten{\underline}{\unsure{\rlap{$G$}F|F}}_p)$, fitting into an exact sequence
\[
0\handwritten{\longrightarrow}\handwritten{\underline}{G}(\handwritten{\underline}{F}_q)\handwritten{\longrightarrow}\handwritten{\underline}{G}(\handwritten{\underline}{F}_p)\handwritten{\longrightarrow}\operatorname{Rep}(\handwritten{\underline}{Z}/\deleted{a}\handwritten{\nu}\handwritten{\underline}{Z})\handwritten{\longrightarrow}0\quad .
\]
It is then clear that for $\handwritten{\underline}{F}_q$ as for $\handwritten{\underline}{F}_p$, the answer to 6.2 is affirmative, hence the consi\rlap{e}derations of number\textsuperscript{o}7 apply: one finds the structure of $\handwritten{\underline}{G}(\handwritten{\underline}{F}_q)$ as a gerbe + its additional structures examined so far.

The \inserted{canonical $(E,F_E)\handwritten{\longmapsto}E$}\handwritten{\begin{tikzpicture}[baseline=0cm,x=1cm,y=1cm]\draw[thin] (0,-.18)..controls (-.08,.1) and (.05,.18)..(.22,.18)..controls (.8,.26) and (1.3,.2)..(3.4,.2)..controls (3.55,.2) and (3.58,.3)..(3.52,.36);\end{tikzpicture}} fibre functor on $\operatorname{Fcriso}(\handwritten{\underline}{F}_q)$ \deleted{with value} over the field $K_{\deleted{a}\handwritten{\nu}}=K(\handwritten{\underline}{\mathbb{F}}_{q=p^{\deleted{a}\handwritten{\nu}}})$ \margin{\handwritten{rather present $\unsure{=}f_q$ as an \uline{automorphism of the identity functor} of $\operatorname{Fcriso}(\mathbb{F}_q)$, gives an element of $G_\nu(\mathbb{Q}_p)$}}
(i.e. canonical section of $\handwritten{\underline}{G}(\handwritten{\underline}{F}_q)$ over $K_{\deleted{a}\handwritten{\nu}}$) is endowed with an additional structure, namely a linear automorphism $F_E{}^a$, so that the \deleted{algebraic} group over \deleted{k\hspace{1em}\unsure{e|c|a}} $K_a$ corresponding to this section of $\handwritten{\underline}{G}(\handwritten{\underline}{F}_q)$ over $K_a$, namely the subgroup \inserted{$\handwritten{G_\nu=\nu}$}\handwritten{\begin{tikzpicture}[baseline=0cm,x=1cm,y=1cm]\draw[thin] (0,-.18)..controls (-.08,.1) and (.05,.18)..(.22,.18)..controls (.8,.26) and (1.3,.2)..(1.3,.2)..controls (1.45,.2) and (1.48,.3)..(1.42,.36);\end{tikzpicture}}$\deleted{a}.G_{\deleted{K_a}\handwritten{K_\nu}}$ of $G_{K_{\deleted{a}\handwritten{\nu}}}$ open of index $\deleted{a}\handwritten{\nu}$ (corresponding to the subgroup of the $\deleted{a}\handwritten{\nu}$.th roots of 1 in the character group $\handwritten{\overline}{\handwritten{\underline}{Q}}^{\handwritten{*}}$ of $G_{\handwritten{K_\nu}}$)
